\section{Projected extraspecial factors: a prime bound and the elliptic escape}
\label{sec:projected-extraspecial}

\paragraph{Scope and review status.}
D-TD05--07 are new partial results for actual factorwise tensor quotients.
The parent suggested the scalar-times-core budget and the full normalizer
trace direction; the proofs below are the author's subsequent derivation.
Complete author and separate root whole-proof readings pass. Root also read the complete source/status record, every checker line, full first output and exact input/exposition binding; no root rerun is claimed.  No bound on the
remaining exponent $m$ at the small primes is asserted.  In particular,
the elliptic examples at the end are projected data, not sixfold sources
with the required full inertia budget or Hodge signature.

\paragraph{Actual source and projective factor.}
Let $C\to\mathbf P^1$ be a connected finite $J$-Galois cover, where
$J\subset\operatorname{GL}(V)$ is its faithful irreducible coefficient image.  Write
$e=\dim V$ and let $T_1,\ldots,T_r$ be a generating product-one inertia
tuple satisfying the sphere-source identity
\[
                 \sum_i\operatorname{rank}(T_i-I)=2e+6.
                 \tag{TDE.1}
\]
Suppose $V=W\otimes U$ is preserved factorwise by $J$, with reciprocal
scalar ambiguity allowed in its two factors.  There is then an actual
finite quotient
\[
        P=\operatorname{im}\bigl(J\longrightarrow\operatorname{PGL}(W)\bigr),
        \qquad Z=C/\ker(J\to P),
        \qquad Z/P=\mathbf P^1.
        \tag{TDE.2}
\]
Assume $P$ contains and normalizes a subgroup
\[
 Q\simeq (\mathbf F_p)^{2m},\qquad m\ge1,
 \tag{TDE.3}
\]
whose projective action on $W$ is irreducible and whose commutator pairing
on scalar lifts is nondegenerate.  Thus, for lifts $L_v,L_w$ of $v,w\in Q$,
\[
 L_vL_wL_v^{-1}L_w^{-1}=\beta(v,w)I,
 \qquad \beta:Q\times Q\longrightarrow\mu_p
 \tag{TDE.4}
\]
is a nondegenerate alternating pairing.  This is the intrinsic
projective Heisenberg hypothesis; no splitting of a normalizer extension
is included in it.  Put
$\nu(A)=\dim W-\max_\lambda\dim\ker(A-\lambda I)$.

This setup applies to a normal extraspecial subgroup of a primitive
actual image $J$.  Indeed, the isotypic summands of its restriction to
$V$ are permuted by $J$, so primitivity makes the restriction homogeneous.
Evaluation gives $V=W\otimes U$ with irreducible $W$ and faithful action
of that subgroup on $W$: an element acting trivially on $W$ would act
trivially on $V$.  Its center is scalar by Schur's lemma.  The algebra
it generates is $\operatorname{End}(W)\otimes I$, so its normalizer acts
factorwise, by the matrix-algebra normalizer argument in D-TD04.
The projective quotient in (TDE.2) is therefore genuine.  Neither
factorwise tensor indecomposability nor an equality of coefficient
fields is being assumed.

\subsection{D-TD05: exact trace of every projective normalizer}

\begin{theorem}[D-TD05]
Under (TDE.3)--(TDE.4), $a=\dim W=p^m$.  If a unitary operator $A$
normalizes $Q$ projectively, let $g$ denote its induced symplectic
automorphism of $Q$.  Then
\[
        |\operatorname{tr}A|^2\in\{0,\ p^{\dim\operatorname{Fix}(g)}\}.
        \tag{TDE.5}
\]
Every nonscalar such $A$, including one with $g=1$, satisfies
\[
                  |\operatorname{tr}A|/a\le p^{-1/2}.
                  \tag{TDE.6}
\]
Moreover, the projective centralizer is exactly
\[
                   C_{\operatorname{PGL}(W)}(Q)=Q.
                   \tag{TDE.7}
\]
Thus the kernel of the normalizer's action on $Q$ is $Q$, not merely
the scalar identity.
\end{theorem}

\begin{proof}
The inverse image of the finite group $P$ in $\operatorname{SL}(W)$
is finite, since its kernel is $\mu_a$.  Averaging a Hermitian form
unitarizes it.  The inverse image $E$ of $Q$ has scalar center $B$;
nondegeneracy of $\beta$ shows that no other element is central.
Its representation on $W$ is irreducible by hypothesis.  If $v\ne0$,
there is $w$ with $\beta(w,v)\ne1$.  Conjugating $L_v$ by $L_w$
therefore proves $\operatorname{tr}L_v=0$.
Character orthogonality for the irreducible representation of $E$
now gives
\[
 1=\frac1{|E|}\sum_{x\in E}|\operatorname{tr}x|^2
   =\frac{|B|a^2}{|E|}=\frac{a^2}{|Q|}.
 \tag{TDE.8}
\]
This proves $a=p^m$.  Choose unitary lifts $L_v$, with $L_0=I$.
They are pairwise orthogonal for the trace inner product, because
$L_vL_w^{-1}$ has zero trace when $v\ne w$.  There are $a^2$ of
them, so they are a basis of $\operatorname{End}(W)$.
The only known representation-theoretic input here is finite character
orthogonality, in \cite[Sections2.3--2.4]{OG:Serre1977}.

Conjugation by $A$ preserves commutators, so $g$ preserves $\beta$.
Write
\[
 A L_v A^{-1}=\eta_v L_{gv},\qquad |\eta_v|=1.
 \tag{TDE.9}
\]
Let $F=\operatorname{Fix}(g)$.  The restriction $\eta|_F$ is an
ordinary character of the additive group $F$.  To see this without
any phase convention, write $L_vL_w=c(v,w)L_{v+w}$.  For $v,w\in F$,
conjugating this equality gives the same factor $c(v,w)$ on both
sides, and hence $\eta_{v+w}=\eta_v\eta_w$.  Also $\eta_0=1$.
In particular this argument works at $p=2$; it neither asserts a
globally linear phase on all of $Q$ nor assumes a quadratic-refinement
convention.

Taking the trace of conjugation on the displayed basis gives
\[
 |\operatorname{tr}A|^2
 =\operatorname{tr}\bigl(\operatorname{Ad}(A)|\operatorname{End}(W)\bigr)
 =\sum_{v\in F}\eta_v.
 \tag{TDE.10}
\]
The last sum is $|F|$ for the trivial character and zero otherwise.
This proves (TDE.5).  If $g\ne1$, then $\dim F\le2m-1$, proving
(TDE.6).  If $g=1$ and $\eta$ is trivial, $A$ centralizes a basis
of the full matrix algebra, so it is scalar.  If $g=1$ and $A$
is nonscalar, $\eta$ is nontrivial and its trace is zero; this is
the other case of (TDE.6).

Finally suppose $[A]$ centralizes $Q$ projectively.  The argument for
$g=1$ makes $\eta$ a character of $Q$.  Nondegeneracy identifies
$Q$ with its character group by $u\mapsto\beta(u,{\cdot})$.
Choose the corresponding $u$.  Then $L_u^{-1}A$ centralizes every
$L_v$ and is scalar.  Conversely every element of $Q$ centralizes
$Q$ projectively, giving (TDE.7).  This last argument also applies
to a nonunitary representative, since its phases form a finite
character and the matrix-algebra conclusion is algebraic.
\end{proof}

\subsection{D-TD06: the actual-cover prime bound}

\begin{theorem}[D-TD06]
Assume (TDE.1)--(TDE.4).  If $g(Z)\ge2$, then
\[
                 e(1-85/\sqrt p)\le252,
                 \qquad p<7744.
                 \tag{TDE.11}
\]
If $m\ge2$, the genus hypothesis is automatic.  Consequently every
source with $p\ge7744$ has $m=1$ and $g(Z)=1$.

In the special scalar-times-core case $P=Q$, the necessary list is
\[
             (m,p)\in\{(1,2),(1,3),(1,5),(1,7),(2,2)\}.
             \tag{TDE.12}
\]
The list is a necessary budget condition, not a statement of existence.
\end{theorem}

\begin{proof}
Remove identity images from the projected inertia tuple of (TDE.2),
and let $o_1,\ldots,o_k\ge2$ be the exact orders of the remaining
images.  These elements generate $P$ and have product one.
For an element of order $o$, choose a unitary representative $A$
with $A^o=I$.  Every $A^j$, $1\le j<o$, is nonscalar.  Fourier
inversion of its eigenvalue multiplicities and (TDE.6) give
\[
 \frac{\nu(A)}a\ge(1-p^{-1/2})(1-1/o).
 \tag{TDE.13}
\]
This estimate applies to an arbitrary eigenvalue, so scalars cannot
affect it.  If $B$ is the operator on $U$, each eigenspace of $A\otimes B$
has dimension at most
$(a-\nu(A))\dim U$: decompose $U$ into its eigenspaces and apply
the largest-eigenspace bound to each summand.  Thus (TDE.13) is a
lower bound for the actual moving rank divided by $e$, regardless
of the tensor partner.  Identity projected images contribute a
nonnegative amount and may be discarded.

Riemann--Hurwitz for the actual cover gives
\[
       2g(Z)-2=|P|\left(\sum_{i=1}^k(1-1/o_i)-2\right).
       \tag{TDE.14}
\]
If $g(Z)\ge2$, the positive quantity in parentheses is at least
$1/42$.  Here is the elementary integer-order check.  For $k\ge5$
it is at least $1/2$.  For $k=4$ the zero case is all orders two,
and the least positive excess is $1/6$.  For $k=3$, order the
orders increasingly: first order at least four gives excess at
least $1/4$; first order three gives least positive case $(3,3,4)$,
of excess $1/12$; first order two leaves $(2,3,w)$, $w\ge7$,
$(2,4,w)$, $w\ge5$, or second order at least five, giving respectively
least excesses $1/42$, $1/20$, and $1/10$.  Fewer marks cannot
give positive excess.  Summing (TDE.13) therefore yields
\[
                  2+6/e\ge (1-p^{-1/2})85/42.
                  \tag{TDE.15}
\]
This is the first inequality of (TDE.11).  If $p\ge7744=88^2$,
then $e\ge p^m\ge p$, and
\[
 e(1-85/\sqrt p)\ge\sqrt p(\sqrt p-85)\ge88\cdot3=264>252,
 \tag{TDE.16}
\]
a contradiction.  No count of twelve nonidentity projected marks
has been used.

We next justify the genus assertions.  An elementary abelian odd
$p$-subgroup of $\operatorname{PGL}_2(\mathbf C)$ has rank at most
one: a nonidentity element has two fixed points, and every commuting
element preserves each of them.  An interchange would conjugate the
element to its inverse, which is different for odd $p$.  Such a group
is contained in the diagonal torus, whose $p$-torsion is cyclic.
For $p=2$, the subgroup fixing both points has order at most two,
and the permutation group on the pair has order at most two; the
elementary abelian rank is at most two.  Thus genus zero is impossible
when $m\ge2$, and is also impossible for $m=1$, $p$ odd.

On a complex genus-one curve choose an origin.  Its automorphism
group is translations semidirect the cyclic origin-fixing group of
order two, four, or six; see
\cite[ChapterIV, Proposition7.13 and Remark7.18, pp.~140--141]{TD:MilneEC2021}.
The kernel of the image of $Q$ in this cyclic group lies in the
$p$-torsion translations, of rank two.  The image has rank at most
one.  The full rank of $Q$ is therefore at most three, and at most
two when $p\ge5$.  The torsion assertion is
\cite[ChapterII, Proposition6.3, p.~70]{TD:MilneEC2021}.
For $m\ge2$ the rank $2m\ge4$ excludes genus one too.  Combining
these facts with (TDE.11) proves the remaining genus claims.

Finally let $P=Q$.  Every nonidentity projected element has order
$p$ and every proper nontrivial power has trace zero, so Fourier
inversion gives every eigenvalue multiplicity exactly $a/p$.
If $k$ nonzero vectors generate $\mathbf F_p^{2m}$ and sum to zero,
then $k\ge2m+1$.  Consequently
\[
      (2m+1)(1-1/p)\le 2+6/e\le 2+6/p^m.
      \tag{TDE.17}
\]
For $m\ge3$, the left side is at least $7/2$ and the right side
at most $11/4$, impossible.  For $m=2$, (TDE.17) is
$3p^2-5p-6\le0$, leaving only the prime two.  For $m=1$ it gives
$p\le9$, leaving precisely the four primes in (TDE.12).
If $J$ is primitive, this scalar-times-core case is impossible
altogether: a maximal isotropic subgroup of $Q$ has $p^m$ distinct
one-dimensional common eigenspaces on $W$ (character Fourier
inversion as above).  Explicitly, commuting lifts of a basis of that
subgroup can be scaled to have order $p$, yielding an ordinary
representation of its additive group.  Every nonzero element has
trace zero, so each character occurs once.  The group $Q$ permutes
these eigenlines.  Tensoring those lines
with $U$ gives a nontrivial system of imprimitivity for $J$.
\end{proof}

\subsection{D-TD07: exact description of the large-prime escape}

\begin{theorem}[D-TD07]
If a source in D-TD06 has $p\ge7744$, then
\[
     m=1,\quad g(Z)=1,\quad
     P=Z[p]\rtimes C_d,\qquad d\in\{2,3,4,6\}.
     \tag{TDE.18}
\]
If the actual factor action $P\subset\operatorname{PGL}(W)$ is
primitive, the only remaining cases are
\[
 \begin{array}{c|c}
 d&\text{necessary and sufficient condition for this factor action
 to be primitive}\\ \hline
 3\text{ or }6&p\equiv2\pmod3\\
 4&p\equiv3\pmod4 .
 \end{array}
 \tag{TDE.19}
\]
In particular the $d=2$ factor is imprimitive.  These conclusions
are necessary restrictions on an actual source.  They do not bound
$m$ for the primes below $7744$.
\end{theorem}

\begin{proof}
The first two assertions follow from D-TD06.  Put
$T=P\cap\operatorname{Trans}(Z)$.  As $p\ge5$, the subgroup
$Q$ acts trivially on the origin-fixing cyclic quotient, so
$Q\subset T$.  Translations commute.  Thus $T$ centralizes $Q$
projectively, and (TDE.7) gives $T=Q$.  Since both $Q$ and $Z[p]$
have order $p^2$, they coincide.  It follows that $P/Q$ is cyclic
of some order $d\in\{1,2,3,4,6\}$.  The case $d=1$ would make
$Z/P$ a quotient by translations and hence another genus-one curve,
contradicting (TDE.2).

Choose an element whose linear part generates the cyclic quotient.
On $Z=\mathbf C/\Lambda$ it has the form $z\mapsto\zeta_d z+b$.
The nonzero endomorphism $1-\zeta_d$ is surjective, so this element
has a fixed point.  Taking it as origin makes the element linear
of exact order $d$.  It supplies the asserted complement.  The
complex-torus description and the characterization of its
origin-preserving maps are
\cite[ChapterIII, Proposition3.3, p.~94, and Theorem3.10, p.~96]{TD:MilneEC2021}.

We check primitivity rather than assuming a normalizer classification.
The action of the complement on $Q=Z[p]$ is the reduction of its
action on the rank-two period lattice.  For orders $2,3,4,6$ the
characteristic polynomials are, respectively,
\[
             (X+1)^2,\quad X^2+X+1,\quad X^2+1,\quad X^2-X+1.
             \tag{TDE.20}
\]
If this action preserves a line $L\subset Q$, the lifts of $L$
have $p$ distinct common eigenlines on $W$: their nontrivial traces
vanish and $\dim W=p$.  These lines are permuted by $P$, making
the factor imprimitive.  Conversely, any nontrivial imprimitivity
system in the prime dimension $p$ consists of $p$ lines.  The
elementary abelian $p$-group $Q$ acts on them.  Its image in $S_p$
is either trivial or cyclic of order $p$.  In the trivial case
all lifts are simultaneously diagonal and commute, contrary to
(TDE.4).  In the other case the kernel is a line in $Q$ and is
$P$-invariant.  Therefore primitivity is equivalent to irreducibility
of the complement action on $Q$.  Reducing the polynomials (TDE.20)
gives exactly (TDE.19).  If the original $J$ is primitive, its
projective factor $P$ is primitive, since any block decomposition
on $W$ tensors with $U$ to a block decomposition on $V$.
\end{proof}

\paragraph{An actual projected family with unbounded primes.}
For every odd prime $p\equiv3\pmod4$ let
\[
       E:\ y^2=x^3-x,\qquad \iota(x,y)=(-x,iy),\qquad
       E\xrightarrow{[p]}E\xrightarrow{x^2}\mathbf P^1.
       \tag{TDE.21}
\]
The composite is a connected Galois cover with group
$E[p]\rtimes\langle\iota\rangle$, of order $4p^2$.
The multiplication map is unramified, and the second map is the
quotient by the order-four automorphism.  Its two order-four fixed
points are $(0,0)$ and infinity; the other nonzero two-torsion
points form one orbit with stabilizer of order two.  Thus the
inertia signature is $(2,4,4)$.

Here is a matching faithful projective Heisenberg representation.
On the basis $e_j$, $j\in\mathbf F_p$, let $\zeta$ be a primitive
$p$th root and set
\[
 Xe_j=e_{j+1},\qquad Ze_j=\zeta^j e_j,\qquad
 Fe_j=p^{-1/2}\sum_k\zeta^{jk}e_k.
 \tag{TDE.22}
\]
Then $ZX=\zeta XZ$, $FXF^{-1}=Z$, $FZF^{-1}=X^{-1}$,
$F^2e_j=e_{-j}$, and $F^4=I$.  These formulas show directly that
$\langle X,Z,F\rangle$ is finite and its projective image is
$\mathbf F_p^2\rtimes C_4$, with complement matrix
$\left(\begin{smallmatrix}0&-1\\1&0\end{smallmatrix}\right)$.
This is the same semidirect product as the cover in (TDE.21),
since the irreducible characteristic polynomial of both actions
is $X^2+1$.  The factor is primitive by the proved criterion.
There are arbitrarily large such primes by Dirichlet's theorem
\cite[Section1, p.~63]{PD:Harrison2009}.

This example establishes that the projected quotient-cover and
projective-factor hypotheses alone allow the genus-one escape
with unbounded $p$.  It does not construct a full $J$-cover and
tensor partner obeying (TDE.1), nor prove a balanced Hodge
signature, period rank nine, coefficient-field descent, or a
prescribed polarization.  Those additional source conditions,
and the exponent $m$ for the remaining small primes, remain gates.
